\section{Results}
\label{sec:results}

\begin{figure}[!t]
\centering
\includegraphics[width=\columnwidth]{fig_trace.pdf}
\caption{A real 8-minute drive of the HCRL KIA Soul (top: decoded vehicle
speed) replayed through a moderately (middle) and severely (bottom) worn
DLC. Losses cluster where the vehicle is moving and vanish at standstill:
the excitation coherence that ECTA measures. Severe wear adds brown-out
silences (tall bars).}
\label{fig:trace}
\end{figure}

Fig.~\ref{fig:trace} illustrates the phenomenon on real data before any
statistics: in the same drive, losses from a worn DLC follow the vehicle's
motion, while the tool sees nothing at standstill.

\subsection{Detection}

\begin{table*}[!t]
\caption{Detection Performance (AUROC) on Real-Vehicle Telemetry. ``All'': DLC Wear (Any Stage) vs. Healthy \emph{and} Confounder Windows. ``Mod.''/``Sev.'': Moderate/Severe Wear vs. Real Healthy Windows Only. Vehicle- or Capture-Disjoint Splits; Best per Column in Bold.}
\label{tab:detection}
\centering
\footnotesize
\setlength{\tabcolsep}{3.6pt}
\resizebox{\textwidth}{!}{%
\begin{tabular}{@{}l ccc ccc ccc ccc@{}}
\toprule
 & \multicolumn{3}{c}{Full-bus passive} & \multicolumn{3}{c}{Sampling passive} & \multicolumn{3}{c}{OBD-II polling} & \multicolumn{3}{c}{VED fleet (381 vehicles)}\\
\cmidrule(lr){2-4}\cmidrule(lr){5-7}\cmidrule(lr){8-10}\cmidrule(l){11-13}
Method & All & Mod. & Sev. & All & Mod. & Sev. & All & Mod. & Sev. & All & Mod. & Sev.\\
\midrule
\tinput{tables/table_detection.tex}
\bottomrule
\end{tabular}}
\end{table*}

Table~\ref{tab:detection} reports detection on all four corpora. Three
findings stand out.

\emph{(i) ECTA is the only method that is strong on both questions.} On the
realistic task---wear against everything else, confounders
included---ECTA has the highest AUROC on the full-bus (0.869), polling
(0.723) and fleet (0.775, 95\% CI 0.772--0.778) corpora and is
statistically tied with the 1D-CNN on the sampling corpus (0.675 vs.\
0.677, overlapping CIs). Single-statistic rules are excellent at telling
severe wear from \emph{healthy} data (e.g., the U-code timeout rule reaches
0.962 on the polling corpus and the loss-rate threshold 0.984 on full-bus
captures), but they collapse to 0.56--0.62 once confounders are present,
because every confounder also produces losses or timeouts. The unsupervised
detectors (Isolation Forest, one-class SVM, LSTM autoencoder) flag
\emph{any} abnormality and therefore also cannot separate wear from its
look-alikes (0.54--0.65 overall).

\emph{(ii) Detectability is governed by the wear stage and by the tool.}
Severe wear (6--30\,$\Omega$, where brown-outs begin) is detected against
healthy telemetry at every fidelity: 0.979 (full-bus), 0.957 (sampling),
0.891 (polling) and 0.851 (VED). Moderate wear (1.5--6\,$\Omega$) is
detected reliably only by the full-bus logger (0.875, CI 0.79--0.95): its
interruptions last tens to hundreds of microseconds and destroy individual
frames, which only a tool that sees every frame can count. To a sampling
logger or a polling tester transacting about once per second they are
almost invisible (0.49--0.56). Incipient wear is indistinguishable from
healthy telemetry in any single window (AUROC 0.50 for all methods; not
tabulated). These limits are physical, not algorithmic: they follow from
the duty cycle of the observation, and no method in Table~\ref{tab:detection}
escapes them.

\emph{(iii) Excitation coherence is what makes polling-mode detection
work.} Removing the coherence features costs 0.037 AUROC on the polling
corpus and 0.071 on VED (on VED the gain over the strongest
alternative, ECTA without structure features, has a 95\% CI of
$+0.017$ to $+0.019$), while on the full-bus corpus, whose
driving segment is short and whose idle and CAN-MIRGU segments carry no
usable excitation variation, it changes nothing---as expected. The
training-free $z_\beta$ test alone is not a detector (AUROC $\approx$0.5):
coherence is a \emph{discriminating} feature, informative only in
combination with evidence that losses exist.

\begin{figure*}[!t]
\centering
\includegraphics[width=\textwidth]{fig_roc.pdf}
\caption{ROC curves (wear vs.\ healthy and confounders) on the four corpora,
pooled out-of-fold scores.}
\label{fig:roc}
\end{figure*}

\subsection{Telling Wear from Its Look-Alikes}

\begin{table}[!t]
\caption{Discrimination of Moderate/Severe DLC Wear From Each Confounder (AUROC)}
\label{tab:conf}
\centering
\footnotesize
\setlength{\tabcolsep}{3pt}
\resizebox{\columnwidth}{!}{%
\begin{tabular}{@{}llccc@{}}
\toprule
Corpus & Method & ECU-side & Bus EMI & Overfl./rand.\\
\midrule
\tinput{tables/table_confounders.tex}
\bottomrule
\end{tabular}}
\end{table}

Table~\ref{tab:conf} isolates the attribution question. Against bus-wide
EMI, coherence features raise the AUROC from 0.741 to 0.878 on the polling
corpus and from 0.737 to 0.947 on VED, even though half of the simulated
EMI was deliberately made ignition-coupled. The separation comes from the
joint footprint: DLC interruptions remove transactions and correlate with
\emph{both} engine and road excitation, whereas EMI delays transactions and
correlates, at most, with engine speed. Against vibration-independent loss,
coherence adds 0.031 on VED. Against an ECU-side intermittent no
telemetry feature helps in polling mode (0.73--0.76 for all methods): an
intermittent engine-ECU connector and an intermittent DLC both make Mode-01
responses disappear, both are vibration-driven, and a tester that talks
only to that ECU cannot tell them apart. On the full-bus corpus, where the
selectivity of losses across identifiers is observable, the same
distinction reaches 0.936. This is a structural limit worth stating
plainly: attributing a fault to the DLC rather than to an ECU connector
requires visibility of traffic from more than one ECU.

\begin{table}[!t]
\caption{Five-Way Fault-Origin Attribution: Macro-F1 (All Windows; Excluding Incipient Windows), Per-Class F1 (Excluding Incipient), and the 1D-CNN Macro-F1}
\label{tab:attr}
\centering
\footnotesize
\setlength{\tabcolsep}{2.4pt}
\resizebox{\columnwidth}{!}{%
\begin{tabular}{@{}lcc ccccc c@{}}
\toprule
 & \multicolumn{2}{c}{ECTA macro-F1} & \multicolumn{5}{c}{ECTA per-class F1} & CNN\\
\cmidrule(lr){2-3}\cmidrule(lr){4-8}
Corpus & all & excl. & Hlth. & DLC & ECU & EMI & Ovf. & all\\
\midrule
\tinput{tables/table_attribution.tex}
\bottomrule
\end{tabular}}
\end{table}

Table~\ref{tab:attr} summarizes the full five-way attribution. ECTA
outperforms the end-to-end 1D-CNN by 0.06--0.40 macro-F1 on three corpora
and is comparable on the sampling corpus, where all methods are limited by
the logger. Fig.~\ref{fig:conf} shows the VED confusion matrix: bus EMI
and overflow/random loss are recognized reliably; the dominant confusions
are healthy~$\leftrightarrow$~early wear (physically indistinguishable in a
single trip) and DLC~$\leftrightarrow$~ECU-side (the structural limit above).

\begin{figure}[!t]
\centering
\includegraphics[width=0.82\columnwidth]{fig_confusion_ved.pdf}
\caption{Row-normalized attribution confusion matrix, VED fleet
(381 vehicles, vehicle-disjoint folds).}
\label{fig:conf}
\end{figure}

\subsection{Pooling Evidence Across Trips}

\begin{figure}[!t]
\centering
\includegraphics[width=\columnwidth]{fig_accumulation.pdf}
\caption{Vehicle-level detection on VED when the per-trip log-odds of $K$
trips of the same vehicle are averaged. Negatives include healthy and all
confounder trips.}
\label{fig:acc}
\end{figure}

A worn DLC stays worn, whereas the background irregularity of healthy
telemetry is trip-local. Fig.~\ref{fig:acc} exploits this. Averaging the
per-trip evidence of $K$ trips of the same vehicle raises severe-wear
detection on VED from 0.907 ($K{=}1$) to 0.987 ($K{=}3$) and 0.9998
($K{=}10$) against all negatives, and from 0.850 to 0.999 against healthy
vehicles alone. Pooling does not, however, create evidence that is not
there: moderate and incipient wear remain at chance against healthy
vehicles for every $K$ (0.46--0.50), confirming that for a
$\sim$1.3\,Hz polling logger their footprint lies below the logger's own
background irregularity (the healthy VED trips contain, on median, 31
gaps longer than 0.8\,s per minute).

\subsection{Fleet Prognosis}

\begin{table}[!t]
\caption{Fleet Prognosis on VED (166 Vehicles, 21\,824 Real Trips). MAE of RUL (Days) for Failing Vehicles; Alarm When Predicted RUL $<30$ Days: Timely ($\le 60$ Days Before Failure), Premature, Missed (\% of Failing Vehicles), Median Lead (Days), False Alarms (\% of Non-Failing Vehicles)}
\label{tab:prog}
\centering
\footnotesize
\setlength{\tabcolsep}{2.2pt}
\resizebox{\columnwidth}{!}{%
\begin{tabular}{@{}lcccccc@{}}
\toprule
Method & MAE & Timely & Prem. & Missed & Lead & FA\\
\midrule
\tinput{tables/table_prognosis.tex}
\bottomrule
\end{tabular}}
\end{table}

\begin{figure*}[!t]
\centering
\includegraphics[width=\textwidth]{fig_prognosis.pdf}
\caption{Fleet prognosis on VED. (a) One vehicle's connector, wearing with
its own recorded usage: true resistance and the per-trip health index
learned from telemetry by a vehicle-disjoint model. (b) Median absolute RUL
error versus true RUL across all failing vehicles (20-$\Omega$ failure
definition; the trend alarm is not a point estimator and is omitted).}
\label{fig:prog}
\end{figure*}

The 166 VED vehicles with at least 60 trips over at least 200 days were
each given a connector that wears with that vehicle's own recorded usage
(21\,824 real trips). The learned health index tracks the true contact
resistance with $R^2=0.86$ (Spearman 0.77) on held-out vehicles
(Fig.~\ref{fig:prog}a), but its information is concentrated above
$\approx$2--4\,$\Omega$; below that it is left-censored, as the detection
results predict.

Table~\ref{tab:prog} reports the consequences honestly. With failure
pre-specified at the onset of the severe stage (6\,$\Omega$), the
detectability floor and the failure threshold nearly coincide, and
\emph{no} telemetry-based method beats the population reliability prior:
the RUL error is 64--71~days for all of them. We then evaluated the
threshold that matters operationally---functional failure of the tool link
at 20\,$\Omega$, where brown-outs dominate---noting that this second
threshold was chosen after seeing the first result. There, a simple alarm
on the trend of the learned health index warns 88\% of failing vehicles
in time (median lead 6.4 days) with a 1.5\% false-alarm rate on vehicles
that do not fail, whereas the fleet-reliability schedule, which uses
the same exposure information but no telemetry, is timely for only 36\%
and false-alarms on 61\%. Point estimates of RUL months ahead, by
contrast, are not improved by telemetry at this tool fidelity
(Fig.~\ref{fig:prog}b): the Bayesian exposure-domain estimator, which
correctly treats sub-floor readings as censored, matches the conditional
reliability model but does not beat it. Long-horizon RUL from telemetry
therefore requires a tool that sees moderate wear, i.e., full-bus capture.

\subsection{Robustness to Mis-Specified Physics}
\label{sec:robust}
\begin{figure}[!t]
\centering
\includegraphics[width=\columnwidth]{fig_robustness.pdf}
\caption{Robustness of VED detection when the test data are generated with
mis-specified physics (training always uses the nominal model of
Table~\ref{tab:params}). Dotted lines: nominal test physics. AUROC against
healthy and confounder trips; held-out half of the vehicles.}
\label{fig:robust}
\end{figure}

The constants of Table~\ref{tab:params} that are not fixed by standards
are unmeasured for J1962 receptacles. We therefore trained ECTA once on
the nominal model and tested it on held-out real data generated with each
constant perturbed (Fig.~\ref{fig:robust}; for the full-bus corpus,
training on three vehicles and testing on the other three). On the
full-bus corpus ECTA is robust to \emph{every} perturbation: the overall
AUROC stays between 0.84 and 0.95 (nominal 0.901) and severe-stage AUROC
never falls below 0.97, because a tool that counts individual lost frames
does not depend on brown-outs. For the polling fleet two regimes emerge.
Detection is insensitive---within $\pm0.02$ AUROC of nominal on VED
(0.748)---to the resistance fluctuation spread $\sigma$,
the opportunity rate $\nu_0$, the magnitude of the unobserved road input,
and the composition of the excitation (engine-only or road-only), i.e., to
the constants that shape \emph{where} and \emph{how often} interruptions
occur. It is sensitive to the constants that decide whether interruptions
cause \emph{brown-outs}: shorter interruptions ($d_0\times\frac14$), a
longer tool hold-up time (\SI{10}{\milli\second}), a higher discontinuity
threshold (\SI{15}{\ohm}) or fewer supply-pin interruptions
($p_{\mathrm{pwr}}=0.1$) lower severe-stage AUROC on VED from 0.887 to
0.62--0.79, while the opposite perturbations raise it to 0.93--0.97. The
practical reading is that a polling tool sees wear chiefly through
brown-outs, so the two quantities most worth measuring first on a bench
are the interruption-duration distribution of worn receptacles and the
hold-up time of the tool itself---the latter a design parameter that,
somewhat counter-intuitively, makes a tool a \emph{better} wear sensor
when it is short.


\subsection{Feature Importance and Cost}
On VED the five most important features by split gain are the dispersion
of timeout-gap durations, the strong-event high/low excitation ratio, the
longest gap, long gaps per minute and the all-event excitation ratio
(Fig.~\ref{fig:imp}); coherence features account for
48
\unskip\% of the total gain. Feature extraction
takes 23~ms for a 30-s full-bus window containing 58\,583 frames and
4.8~ms for a 120-s polling window in unoptimized Python, and the trained
detector is a 1.3-MB tree ensemble; ECTA is therefore deployable on the
tool or at the fleet back end at negligible cost.

\begin{figure}[!t]
\centering
\includegraphics[width=\columnwidth]{fig_importance_ved.pdf}
\caption{Twelve most important ECTA features on VED (share of split
gain). Blue: excitation-coherence features.}
\label{fig:imp}
\end{figure}
